\subsubsection{The diagonal-invariant correspondence: an explicit torsion test}

The preceding bounded-complex theorem closes the ordinary type-A
Soergel case. A separate calculation explains why its most natural
coupled affine correspondence already fails. It identifies the
actual singular fiber and its torsion, rather than inferring failure
only from the absence of a reflection.

Let \(k\) have characteristic zero and put
\[
 R=k[x,y,z_1,\ldots,z_m],\qquad \deg x=\deg y=\deg z_j=2.
\]
Let \(\tau\) negate \(x,y\) and fix every \(z_j\). These are the two
independent anti-invariant coordinates of \((s,s)\); write
\(Q=\operatorname{Frac}R\). Define
\begin{equation}\label{eq:diagonalinvariantcandidate}
 S=R^\tau=k[z_1,\ldots,z_m,x^2,xy,y^2],
 \qquad T=R\otimes_S R,
 \qquad \mathscr D=R(2)\oplus R^{\oplus2}\oplus R(-2).
\end{equation}
The invariants have the displayed generators because a monomial
in \(x,y\) is invariant precisely when its total \(x,y\)-degree is
even. In the coordinates \(a=x^2,b=xy,c=y^2\), their relation is
\(ac=b^2\).

The correspondence has the required two generic graphs:
\[
                    T_Q\simeq Q_e\oplus Q_\tau.
\]
Indeed \(\operatorname{Frac}S=Q^\tau\), and \(Q\) is its separable
quadratic extension generated by \(x\), with
\(y=(xy/x^2)x\). The elementary quadratic Galois splitting of
\(Q\otimes_{Q^\tau}Q\) gives the identity and \(\tau\) graphs.
Thus this is an actual candidate for the diagonal support.

\begin{finding}[\fid{F59}, explicit invariant-correspondence test]
\label{f:invariantcorrespondencetest}
The left \(R\)-module \(T\) has generic rank two and a
three-dimensional fiber at the homogeneous origin. It has
nonzero torsion supported on \(x=y=0\), described by the exact
sequence below. For every nonzero finite internally graded
\(k\)-space \(U\), the degree-zero bimodule
\(X_U=U\otimes_kT\) fails the derived quadratic
\begin{equation}\label{eq:invariantderivedquadratic}
 X_U\otimes_R^L X_U\simeq\mathscr D\otimes_R X_U.
\end{equation}
The same exclusion holds with \(U\) replaced by any bounded
complex of finite graded \(k\)-spaces with nonzero cohomology.
\end{finding}

\begin{proof}
\par\smallskip\noindent\textbf{The nonflat fiber.}\quad
At the invariant origin
\[
 \mathfrak n=(z_1,\ldots,z_m,x^2,xy,y^2)\subset S
\]
one has
\begin{equation}\label{eq:invariantoriginfiber}
 R/\mathfrak nR
       =k[x,y]/(x^2,xy,y^2),
 \qquad \dim_k(R/\mathfrak nR)=3.
\end{equation}
The generic \(S\)-rank of \(R\) is two, so this fiber rules out
local freeness of rank two. Tensoring \(T\) with the residue
field at the left origin gives the same three-dimensional
space. Its generic left rank is two, so \(T\) is not left-free.

\par\smallskip\noindent\textbf{The exact torsion module.}\quad
Decompose \(R=S\oplus M\) into its even and odd parts under
\(\tau\). The odd part \(M\) is generated over \(S\) by \(x,y\);
its relations are generated by
\begin{equation}\label{eq:oddinvariantrelations}
                 (xy,-x^2),\qquad(y^2,-xy).
\end{equation}
For completeness, if \(p,q\in S\) satisfy \(px+qy=0\),
coprimality gives \(p=yh,q=-xh\) for \(h\in R\).
The parity of \(p,q\) makes \(h\) odd. Write
\(h=xu+yw\) with \(u,w\in S\), which expresses the relation
as the corresponding combination of
\eqref{eq:oddinvariantrelations}. This proves that the
presentation is complete.

Right exactness of tensoring this presentation with the left
copy of \(R\) gives
\begin{equation}\label{eq:invarianttorsionpresentation}
 T\simeq R\oplus
       \frac{Ru\oplus Rv}{(x\omega,y\omega)},
 \qquad \omega=yu-xv.
\end{equation}
Here \(u,v\) represent the right odd generators \(x,y\), each
in internal degree two; \(\omega\) has internal degree four.
In the free module \(Ru\oplus Rv\), its nonzero vector
\(\omega\) has no scalar annihilator. Consequently its image
in the quotient is nonzero, since \(1\notin(x,y)\), and its
annihilator is exactly \((x,y)\).

The map \(u\mapsto x,v\mapsto y\) has kernel \(R\omega\)
before imposing the relations, again by coprimality.
It therefore induces the exact sequence
\begin{equation}\label{eq:invarianttorsionsequence}
 0\longrightarrow (R/(x,y))\omega
 \longrightarrow \frac{Ru\oplus Rv}{(x\omega,y\omega)}
 \longrightarrow (x,y)\longrightarrow0.
\end{equation}
This proves the claimed codimension-two torsion directly.

\par\smallskip\noindent\textbf{The derived quadratic fails.}\quad
Put \(r=\dim_kU>0\). If
\eqref{eq:invariantderivedquadratic} held, taking degree-zero
cohomology would give an ordinary module isomorphism
\begin{equation}\label{eq:invariantordinaryquadratic}
 X_U\otimes_R X_U\simeq\mathscr D\otimes_R X_U.
\end{equation}
Here we use
\(H^0(M\otimes_R^L N)=M\otimes_RN\) for modules in degree
zero; we have not assumed \(T\) flat or discarded higher Tor
groups.

Let \(k\) now also denote the left-origin residue module and
put
\[
 B=k\otimes_S R=k[x,y]/(x^2,xy,y^2),\qquad \dim_kB=3.
\]
Associativity of ordinary tensor products gives
\begin{equation}\label{eq:invariantdoublefiber}
 k\otimes_R(T\otimes_R T)
      \simeq k\otimes_S R\otimes_S R
      \simeq B\otimes_k B.
\end{equation}
For the last equality, the \(S\)-action on \(B\) factors
through \(S/\mathfrak n=k\). After taking this left fiber
and forgetting the internal grading, the two sides of
\eqref{eq:invariantordinaryquadratic} have respective
dimensions
\[
                            9r^2,\qquad 12r.
\]
Their equality would force \(r=4/3\), impossible for a
positive integer. This proves the degree-zero case.

Finally let \(U^\bullet\) be a bounded coefficient complex
with nonzero cohomology. After generic localization,
\(T_Q=Q_e\oplus Q_\tau\). The nonzero semisimple generic
graph complex \(U^\bullet\otimes_k T_Q\), if it satisfied
the quadratic, would have both extreme cohomological
degrees zero: tensoring doubles either extreme, whereas
\(\mathscr D\) is in cohomological degree zero.
Thus \(U^\bullet\) has only \(H^0\). Over \(k\) it is
quasi-isomorphic to that finite graded space, so the
already proved case applies.
\end{proof}

For a lift of the actual class \(P_s\), generic support
would require two copies of each graph, so the candidate
would need \(r=2\). Its origin calculation is then the
concrete discrepancy \(36\ne24\).
This is a direct obstruction to this nonflat invariant
correspondence and its bounded scalar coefficient
complexes. It does not assert that every coherent object
on the union of the two graphs is a scalar multiple of
\(T\).

The stronger finding \ref{f:boundedSoergelNoGo} has already
excluded the remaining bounded ordinary Soergel case,
using the finite Frobenius fiber to prove cohomological
concentration and freeness. The present computation adds
the explicit singular correspondence and torsion
mechanism. It leaves different monoidal categories,
altered coefficient complexes and unbounded constructions
subject to their own hypotheses.

% Separate paragraph for Part III.
\paragraph{The diagonal-invariant test in Grothendieck's language.}
The bimodule \(T\) is the affine structure sheaf of
\[
 \operatorname{Spec}R
       \times_{\operatorname{Spec}S}
 \operatorname{Spec}R.
\]
The two generic graphs do not make its first projection
finite flat of degree two. Equations
\eqref{eq:invariantoriginfiber} and
\eqref{eq:invarianttorsionsequence} give the actual fiber
and torsion defect. The affine module/sheaf dictionary is
EGA~I, Theorem~1.3.7, p.~86, Corollary~1.3.8, p.~87, and
Theorem~1.4.1, pp.~90--92 \cite{EGA}.
Composition of affine correspondences is tested by the
displayed bimodule tensor products; derived composition
retains higher Tor and is not silently replaced by a
flat tensor product. The degree-zero consequence is
proved separately before taking its ordinary fiber.

The coefficient gradings use SGA~3, I, 4.7.2--3 and
cote~90, scan~17, and the effective/virtual vocabulary of
SGA~6, V, \S\S2--3 and cote~133, scan~18
\cite{SGA3,SGA6,Archive}. The bounded-complex framework is
SGA~6, I, Definition~4.7, pp.~121--122, and
Example~4.8, p.~122, with the tensor base specified.
These anchors do not assert freeness of this singular
correspondence or supply its quadratic. Its failure is
proved by the explicit module and fiber calculation.

